% Generated draft: requires mathematical and source review.
\documentclass[11pt]{article}
\usepackage[margin=25mm]{geometry}
\usepackage{fontspec}
\setmainfont{Latin Modern Roman}
\usepackage{amsmath,amssymb,mathtools,unicode-math}
\setmathfont{Latin Modern Math}
\usepackage{enumitem,needspace,xurl,hyperref}
\setlength{\parindent}{0pt}
\setlength{\parskip}{5pt}
\setlength{\emergencystretch}{3em}
\newcommand{\N}{\mathbb{N}}
\newcommand{\Z}{\mathbb{Z}}
\newcommand{\Q}{\mathbb{Q}}
\newcommand{\R}{\mathbb{R}}
\newcommand{\C}{\mathbb{C}}
\newcommand{\F}{\mathbb{F}}
\newcommand{\A}{\mathbb{A}}
\newcommand{\PP}{\mathbb{P}}
\newcommand{\E}{\mathbb{E}}
\newcommand{\eps}{\varepsilon}
\newcommand{\dd}{\,\mathrm{d}}
\begin{document}
\section*{UnsolvedMath ID: 30000119}
\textbf{Source problem code:} catalogue entry
\subsection{Definitions and mathematical statement}
Let \(S\) be the complex Beauville surface intended in the source, and let \(\overline S\) be its complex conjugate. Beauville surfaces are rigid surfaces isogenous to a product: they have a presentation as a free finite-group quotient of a product of smooth projective curves of genera at least two, with additional Beauville conditions. In the unmixed setting one commonly writes \(S=(C_1\times C_2)/G\), with the group preserving the two factors and each quotient map \(C_i\to C_i/G\cong\PP^1\) branched over three points. The extracted record gives no such presentation, group, branch data, or indication whether mixed actions are included.

The two requested determinations are distinct:
\[
\text{Does there exist a biholomorphism }f:S\longrightarrow\overline S?
\]
\[
\text{Does there exist an antiholomorphic involution }\tau:S\longrightarrow S,
\qquad\tau^2=\operatorname{id}_S?
\]
In the second question, \(S\) is called real if it has a model over \(\R\), equivalently a real structure given by such an involution. A biholomorphism to \(\overline S\) corresponds to an antiholomorphic automorphism of \(S\), which is not required to have order two. Descent to \(\R\) requires the existence of a choice of \(f\) obeying the cocycle condition
\[
\overline f\circ f=\operatorname{id}_S.
\]
Consequently the first existence condition must not be identified with the second.

The source asks these questions about ``the Beauville surface \(S\)'' without specifying it. It may seek criteria in terms of construction data, or a classification over a stated family, rather than assert a universal yes answer. The supplied assessment checked 2026-08-21 reports examples and group-theoretic criteria giving both conjugation-invariant and non-conjugation-invariant Beauville surfaces, including non-real examples, and says that a complete classification is broader than those results. No particular example or complete criterion is reproduced in the record.
\subsection{Short English statement}
Determine whether the intended Beauville surface is isomorphic to its conjugate and whether it admits a real structure. The latter requires an antiholomorphic involution, a stronger condition than merely having an antiholomorphic automorphism.
\subsection{Sources}
\begin{itemize}
\item Local UnsolvedMath registry, permanent ID 30000119, catalogue entry, \emph{Reality of Beauville Surfaces}.
\item \emph{Komplexe Analysis} (2005), Oberwolfach Reports, pp. 2171--2216, DOI \url{https://doi.org/10.4171/owr/2004/42}.
\item Bauer--Catanese--Grunewald, supplied preprint link \url{https://arxiv.org/abs/math/0408025}; no full article title or theorem number is reproduced.
\item Supplied literature triage checked 2026-08-21, reporting examples and group-theoretic criteria but not a complete classification.
\end{itemize}
\end{document}
